% !TEX root = ../main.tex
\chapter{Estado del Arte}
\label{chap:estado_del_arte}

Este capítulo presenta una revisión de la literatura y las bases teóricas fundamentales sobre las que se apoya el sistema desarrollado. Se exploran los paradigmas modernos de previsión de demanda, la evolución de los modelos basados en árboles, la cuantificación de la incertidumbre y la optimización de inventarios bajo condiciones de censura.

\section{Previsión de la Demanda en Retail}
Tradicionalmente, la previsión de series temporales se ha abordado mediante modelos \textbf{locales}, donde se entrena un modelo individual para cada SKU o combinación tienda-producto. Los enfoques clásicos como ARIMA \cite{hyndman2021forecasting} o el suavizado exponencial (ETS) han sido el estándar durante décadas gracias a su interpretabilidad y su sólida base estadística. Sin embargo, en el contexto del \textit{retail} masivo presentan tres limitaciones estructurales: (1) su coste computacional crece linealmente con el número de series, haciendo inviable el ajuste individual sobre catálogos de decenas de miles de SKUs; (2) requieren un historial mínimo por serie para estimar sus parámetros, fallando sistemáticamente en productos nuevos o con alta intermitencia (\textit{cold-start problem}); y (3) no pueden aprovechar información compartida entre series como la estacionalidad semanal o la elasticidad de precio, que son patrones transversales a todo el catálogo.

El paradigma actual ha girado hacia los \textbf{Modelos Globales de Pronóstico (GFMs)}, donde un único modelo se entrena sobre la totalidad de las series temporales de la base de datos. Montero-Manso y Hyndman \cite{montero2021principles} demostraron teóricamente que los modelos globales no son restrictivos: bajo condiciones razonables pueden aproximar cualquier función local, pero además se benefician del aprendizaje de representaciones compartidas. Este enfoque ha dominado competiciones de referencia como la M5 Walmart \cite{makridakis2022m5}, donde las soluciones ganadoras fueron casi exclusivamente modelos globales basados en \textit{Machine Learning}.

Una alternativa relevante a los modelos de \textit{boosting} son las arquitecturas de \textit{deep learning} para series temporales, como DeepAR \cite{salinas2020deepar} o el Temporal Fusion Transformer (TFT) \cite{lim2021temporal}. Aunque estas arquitecturas pueden capturar dependencias temporales a largo plazo de forma automática, presentan desventajas en el contexto de este trabajo: requieren volúmenes de datos superiores para converger, son difíciles de interpretar para los responsables logísticos y su tiempo de entrenamiento es sensiblemente mayor sobre datasets de 4,5 millones de observaciones con covariables heterogéneas. En la competición M5 \cite{makridakis2022m5}, los modelos de \textit{boosting} superaron consistentemente a las redes neuronales en términos de precisión y eficiencia, lo que justifica su elección en este trabajo.

\section{Modelos de Gradient Boosting: Comparativa y Selección}
Los algoritmos de \textit{Gradient Boosted Decision Trees} (GBDT) son el estado del arte para datos tabulares y series temporales con covariables exógenas. Las tres implementaciones principales (XGBoost \cite{chen2016xgboost}, LightGBM \cite{ke2017lightgbm} y CatBoost \cite{prokhorenkova2018catboost}) dominaron la competición M5 Walmart y comparten el núcleo algorítmico de boosting por gradiente, pero difieren en aspectos críticos para el contexto de este trabajo.

La Tabla~\ref{tab:gbdt_comparison} resume las dimensiones más relevantes para el problema de previsión de demanda retail con 50.000 series heterogéneas y variables categóricas de alta cardinalidad.

\begin{table}[htbp]
    \centering
    \small
    \setlength{\tabcolsep}{5pt}
    \begin{tabular}{lccc}
        \toprule
        \textbf{Criterio} & \textbf{XGBoost} & \textbf{LightGBM} & \textbf{CatBoost} \\
        \midrule
        Crecimiento de árbol & Por niveles & Por hojas & Oblivious (simétrico) \\
        Variables categóricas & Manual (OHE) & Nativa (particiones) & Nativa (Ordered TS) \\
        Riesgo de \textit{target leakage} & Alto & Medio & Bajo (Ordered Boosting) \\
        Velocidad en datasets grandes & Alta & Muy alta & Media \\
        Regularización implícita & Moderada & Moderada & Alta \\
        Predicción de cuantiles & Sí & Sí & Sí \\
        Incertidumbre nativa (varianza) & No & No & Sí \\
        \bottomrule
    \end{tabular}
    \caption{Comparativa de implementaciones GBDT para el contexto de retail demand fore\-cast\-ing.}
    \label{tab:gbdt_comparison}
\end{table}

La tabla recoge las tres implementaciones que la literatura considera de referencia, pero el
sistema desarrollado implementa \textbf{dos}: LightGBM y CatBoost, ambas bajo el mismo protocolo
de backtesting y el mismo envoltorio cuantílico. XGBoost se descarta antes de implementarse, y por
la fila que más pesa en este dominio: es la única de las tres que exige codificar manualmente unos
identificadores de tienda, producto y categoría de cardinalidad alta, y el \textit{one-hot} de ese
catálogo produce una matriz dispersa que ni escala ni permite al árbol agrupar categorías con
comportamiento parecido.

Entre las dos que sí se implementan, la elección del campeón \textbf{no se decide con los
argumentos de esta tabla}, sino empíricamente y por el criterio que la operación percibe: el coste
logístico simulado sobre el protocolo de backtesting, mediante la lógica de promoción descrita en
la Sección~\ref{sec:promocion}. Esa distinción importa, porque el orden por error puntual y el
orden por coste no coinciden: reportarlos por separado es una de las conclusiones del trabajo. El
campeón resultante es \textbf{CatBoost}, y las propiedades de la tabla explican \emph{por qué es
razonable} que gane, no la razón por la que se selecciona:

\begin{itemize}
    \item \textbf{Gestión nativa de categóricas con estadísticos ordenados:} las dos manejan las
    categóricas sin codificación manual, pero lo hacen de forma distinta. LightGBM ordena las
    categorías por el estadístico del gradiente y parte sobre ese orden, lo que en categorías con
    pocas observaciones puede sobreajustar; CatBoost calcula estadísticos de destino sobre una
    permutación previa (\textit{ordered target statistics}), construcción diseñada precisamente
    para que el valor de la etiqueta de una observación no entre en el estadístico que la codifica.
    Con identificadores de la cardinalidad de este catálogo, esa diferencia es la más relevante de
    las tres.
    \item \textbf{Regularización implícita mediante árboles \textit{oblivious}:} los árboles
    simétricos de CatBoost aplican la misma condición de división en todo un nivel, lo que actúa
    como regularizador estructural sobre un panel donde muchas series tienen historia corta.
\end{itemize}

La estimación nativa de incertidumbre que CatBoost ofrece (última fila de la tabla) \textbf{no se
explota} en este sistema: la cuantificación de la incertidumbre se obtiene de la rejilla de
cuantiles y de la calibración conformal, que son comunes a los dos \textit{backends} y por tanto
mantienen la comparación entre ellos en igualdad de condiciones. Se recoge en la tabla como
propiedad de la herramienta, no como capacidad utilizada.

\section{Cuantificación de la Incertidumbre y Conformal Prediction}
Una predicción puntual es insuficiente para la gestión de inventarios; es necesario cuantificar la probabilidad de diferentes escenarios de demanda para poder aplicar el criterio del Newsvendor. Los enfoques clásicos de cuantificación de incertidumbre se dividen en tres familias: métodos bayesianos, regresión cuantílica y métodos de calibración post-hoc. Los métodos bayesianos \cite{hyndman2021forecasting} ofrecen intervalos bien fundamentados teóricamente pero requieren especificar distribuciones a priori y son computacionalmente costosos a escala. La regresión cuantílica \cite{koenker2005quantile} estima directamente los cuantiles de la distribución condicionada, pero no garantiza cobertura calibrada: el cuantil $q_{0{,}9}$ estimado puede contener la demanda real solo el 70\% de las veces si la distribución de test difiere de la de entrenamiento.

La \textbf{Conformal Prediction (CP)} \cite{angelopoulos2021gentle, vovk2005algorithmic} resuelve este problema ofreciendo garantías de cobertura \textit{distribution-free}: sin asumir ninguna forma funcional de la distribución de la demanda, garantiza que el intervalo predicho al nivel $1-\alpha$ contenga el valor real al menos con probabilidad $1-\alpha$ sobre cualquier muestra intercambiable. Existen varias variantes en la literatura: la CP completa (\textit{full conformal}) es computacionalmente inviable para modelos grandes; la CP dividida (\textit{split conformal}) calibra sobre un conjunto de holdout separado y es eficiente pero desperdicia datos; la CP cruzada (\textit{cross-conformal}) combina las ventajas de ambas mediante validación cruzada.

Un supuesto crítico atraviesa toda la familia: la \textbf{intercambiabilidad} entre los datos de calibración y los de test, que las series temporales violan por construcción. Barber et al.\ \cite{barber2023beyond} caracterizan teóricamente esta situación y acotan la pérdida de cobertura en función de una medida de desviación respecto a la intercambiabilidad, estableciendo que la garantía no desaparece sino que se degrada de forma controlada. Sobre esa base, los métodos de CP adaptativa relajan el supuesto en la práctica: \textit{EnbPI} \cite{xu2021conformal} (Ensemble Batch Prediction Intervals) ajusta dinámicamente el ancho del intervalo basándose en los errores recientes observados, ACI \cite{gibbs2021adaptive} actualiza el nivel de significancia en línea según los fallos de cobertura acumulados, y Zaffran et al.\ \cite{zaffran2022adaptive} comparan sistemáticamente estas variantes sobre series de demanda energética, un dominio con la misma combinación de estacionalidad fuerte y horizonte corto que caracteriza al retail de frescos. Estos métodos son por tanto especialmente robustos ante la no-estacionariedad propia de series de demanda retail con estacionalidad y eventos promocionales; como se verá en el Capítulo~\ref{chap:resultados}, esta limitación es precisamente la que acota la cobertura alcanzada por el sistema desarrollado.

Sin embargo, la CP estándar solo garantiza cobertura \textit{marginal}: el intervalo puede ser correcto en media pero fallar sistemáticamente en subgrupos específicos, como las categorías de alta estacionalidad o los productos intermitentes. La \textbf{Mondrian Conformal Prediction} \cite{vovk2005algorithmic} extiende la garantía a cobertura \textit{condicional} al estratificar la calibración por grupos homogéneos, asegurando que cada subgrupo tenga su propia garantía de cobertura independiente. Esta propiedad es crítica en un catálogo retail heterogéneo.

\section{Optimización de Inventarios: El Problema del Newsvendor}
\label{sec:newsvendor}
La conexión entre el pronóstico y la acción operativa se realiza mediante la teoría de gestión de inventarios estocásticos. El modelo del \textbf{Newsvendor} \cite{silver1998inventory} es el marco de referencia para decisiones de pedido bajo demanda incierta en un único periodo: determina la cantidad óptima a pedir ($Q^*$) minimizando el coste esperado, que equilibra el coste de sobrestock ($c_o$, exceso de producto que se desperdicia o deprecia) y el coste de infra-stock ($c_u$, venta perdida y pérdida de fidelidad del cliente).

La solución óptima viene dada por el \textit{fractil crítico}:
\begin{equation}
    P(D \leq Q^*) = \frac{c_u}{c_u + c_o}
\end{equation}
donde $D$ es la demanda aleatoria. Esta fórmula tiene una interpretación directa: si el coste de rotura es cuatro veces mayor que el de sobrestock ($c_u = 4, c_o = 1$), el fractil óptimo es 0,8, es decir, hay que pedir la cantidad que cubre el 80\% de los escenarios de demanda posibles. Nótese que el fractil óptimo es exactamente el cuantil $q_{0{,}8}$ de la distribución predicha, lo que establece un puente directo entre los intervalos de confianza generados por el modelo probabilístico y la decisión de pedido.

Para entornos de múltiples periodos, el modelo del Newsvendor se extiende a políticas de reposición dinámica. La política $(s, S)$ \cite{silver1998inventory} es la más extendida en la práctica: cuando el nivel de inventario cae por debajo del punto de reorden $s$, se emite un pedido hasta el nivel objetivo $S$. Esta política es óptima bajo costes de pedido fijos y demanda estocástica estacionaria. Su caso límite cuando el coste fijo de pedido es nulo es la política base-stock u Order-Up-To. Este trabajo no modela el estado de inventario entre periodos: resuelve una decisión de un solo periodo por origen de predicción, suficiente para ordenar políticas por su coste y su nivel de servicio, pero no para reproducir la dinámica de reposición que estas familias describen. La extensión a un régimen con estado se recoge como línea de trabajo futuro.

Esta conexión entre predicción y decisión ha cristalizado en la última década en la línea de investigación conocida como \textit{decision-focused learning} o \textit{predict-then-optimize}. Bertsimas y Kallus \cite{bertsimas2020prescriptive} formalizan el paso de la analítica predictiva a la prescriptiva, mostrando que el objetivo relevante no es el error del pronóstico sino el coste de la decisión que induce, y Elmachtoub y Grigas \cite{elmachtoub2022spo} proponen entrenar el modelo directamente contra la pérdida del problema de optimización aguas abajo (\textit{SPO loss}) en lugar de contra una métrica estadística intermedia. El enfoque de este trabajo se sitúa en esta familia pero por una vía más ligera y directamente interpretable: en lugar de diferenciar a través del solver, explota la equivalencia analítica exacta entre el fractil crítico del Newsvendor y el cuantil que minimiza la \textit{pinball loss}, lo que permite alinear entrenamiento y decisión sin abandonar los modelos de \textit{boosting} estándar.

Sin embargo, tanto el Newsvendor como la política $(s, S)$ asumen implícitamente que la distribución $D$ es conocida y libre de sesgo. Esta hipótesis se viola cuando los datos de entrenamiento están censurados por roturas de stock, lo que desplaza el fractil crítico hacia abajo y genera pedidos insuficientes de forma estructural. La sección siguiente analiza este fenómeno y su impacto en la cadena de decisión.

\section{Demanda Censurada y el Efecto de Stockout}
Un desafío crítico y frecuentemente ignorado en \textit{retail} es que las ventas observadas solo coinciden con la demanda real cuando hay stock disponible. Cuando ocurre una rotura de stock (\textit{stockout}), los datos se ven \textbf{censurados por la derecha}: la venta registrada no es la demanda, sino una \emph{cota inferior} de ella, ya que observamos cero (o muy pocas) ventas mientras la demanda latente (lo que los clientes habrían comprado) es estrictamente mayor. Formalmente, al observar $\min(D, \text{stock})$ el valor verdadero queda siempre por encima del dato registrado, que es la definición de censura por la derecha en el sentido del análisis de supervivencia. El estudio teórico de los sistemas de inventario con demanda inobservable se remonta a trabajos fundamentales como el de Nahmias \cite{nahmias1994demand}, y posteriormente, Musalem et al.\ \cite{musalem2010structural} estimaron empíricamente que ignorar este fenómeno subestima la demanda real entre un 5\% y un 30\% dependiendo de la tasa de rotura histórica del producto.

Si un modelo se entrena con estos datos sesgados sin corrección, cae en una ``espiral de muerte'' (\textit{death spiral}): el modelo predice baja demanda, se pide poco producto, se agota antes, y el modelo vuelve a reducir su previsión. Este ciclo se retroalimenta hasta colapsar la reposición de ciertos SKUs. La literatura propone varias estrategias de corrección con distintos compromisos entre precisión y coste computacional:

\begin{itemize}
    \item \textbf{Modelos Tobit}: modelan explícitamente la censura como una distribución truncada, recuperando la media de la distribución latente. Precisos pero requieren especificar la forma distribucional y son computacionalmente costosos a escala.
    \item \textbf{Algoritmos EM (Expectation-Maximization)}: iteran entre estimar la demanda latente y reajustar los parámetros del modelo hasta convergencia. Más flexibles que Tobit pero igualmente costosos para catálogos de 50.000 series.
    \item \textbf{Escalado proporcional}: estima la demanda latente multiplicando la demanda observada por un factor inversamente proporcional a las horas de disponibilidad. Simple y escalable, aunque asume linealidad en la relación entre disponibilidad y demanda.
    \item \textbf{Modelo supervisor}: entrena un modelo auxiliar sobre los días sin stockout para predecir la demanda latente en los días afectados. Captura patrones no lineales y es el enfoque adoptado en este trabajo como estrategia principal.
\end{itemize}

La censura introduce además un segundo problema para la cuantificación de la incertidumbre: los métodos de regresión cuantílica estándar aprenden distribuciones sesgadas si los datos de calibración contienen observaciones censuradas, produciendo intervalos de confianza demasiado estrechos precisamente en los SKUs con mayor historial de roturas. La \textbf{Mondrian Conformal Prediction} mitiga este efecto al calibrar los intervalos de forma separada por categoría de producto: incluso si la distribución global está sesgada, la calibración estratificada corrige el sesgo de cobertura dentro de cada grupo homogéneo, garantizando que el fractil crítico del Newsvendor se calcule sobre una distribución correctamente estimada.

\section{Síntesis y Posicionamiento del Trabajo}
\label{sec:posicionamiento}

La revisión de la literatura revela que los avances en cada uno de los pilares descritos han evolucionado de forma mayoritariamente independiente. Los trabajos sobre modelos globales de \textit{boosting} \cite{makridakis2022m5, ke2017lightgbm, prokhorenkova2018catboost} optimizan la precisión predictiva sin conectarla con la decisión de inventario. Los trabajos sobre Conformal Prediction \cite{angelopoulos2021gentle, xu2021conformal} garantizan cobertura estadística pero asumen datos no censurados y no evalúan el impacto logístico de los intervalos. Los trabajos sobre optimización de inventarios bajo incertidumbre \cite{silver1998inventory, ban2019big} asumen distribuciones de demanda conocidas y libres de sesgo. La literatura de \textit{decision-focused learning} \cite{bertsimas2020prescriptive, elmachtoub2022spo} sí conecta explícitamente predicción y decisión, y es la más próxima al planteamiento de este trabajo, pero lo hace sobre datos limpios y sin cuantificar la incertidumbre con garantías formales ni evaluar la decisión en un régimen dinámico con estado de inventario. Finalmente, la literatura sobre demanda censurada \cite{musalem2010structural} propone correcciones estadísticas pero raramente las integra en un sistema de previsión probabilística con evaluación end-to-end.

Ningún trabajo previo identificado integra simultáneamente estos cuatro elementos: (1) corrección de demanda censurada, (2) forecasting probabilístico global con garantías de cobertura condicional, (3) transformación de la distribución predicha en una decisión de inventario vía el fractil crítico del Newsvendor, y (4) evaluación del impacto económico de esa decisión en las dos magnitudes que la operación percibe, coste logístico y nivel de servicio. Esta integración constituye la contribución central del presente trabajo y justifica el diseño del sistema descrito en los capítulos siguientes.
